\section{Positive and finite center bounds}\label{sec:anchors}

Before comparing different sizes, we place the centers in a fixed
positive interval below the cap. At high density, an assignment first
moment makes satisfiability unlikely. At a sufficiently small positive
density, a matching from clauses to variables supplies a satisfying
assignment regardless of the signs.

The matching criterion is Hall's theorem~\cite{Hall}. Its application
to satisfiability and the associated deficiency counting appear in
Franco and Van Gelder~\cite[Section~8]{FV}. We include the
alternating-path argument and the probability estimate.

\begin{lemma}[Bounds on the centers]\label{lem:center-bounds}
For each fixed \(k\ge3\), there is \(a_0>0\) such that, for all
sufficiently large \(n\),
\[
 a_0\le\mu_n\le2^k+1.
\]
\end{lemma}

\begin{proof}
For any fixed assignment, each proper uniformly signed \(k\)-clause
is satisfied with probability \(1-2^{-k}\). The clauses are independent,
so at \(m=2^kn\) the expected number of satisfying assignments is
\[
 2^n(1-2^{-k})^{2^kn}\le(2/e)^n=o(1).
\]
Consequently \(P_n(2^kn)=o(1)\). Since \(V_n\le L_kn\),
\[
 \mu_n\le2^k+(L_k-2^k)P_n(2^kn)=2^k+o(1).
\]

For the lower bound, fix \(c_0>0\) with \(e^2c_0<1\) and put
\(m=\lfloor c_0n\rfloor<n\). A matching assigning each clause position
a distinct variable occurring in that clause guarantees satisfiability:
give each matched variable the value that satisfies its literal in the
matched clause. These choices are compatible because the matched
variables are distinct.

If no such matching exists, some \(t\) clause positions use fewer than
\(t\) variables. To see this directly, take a maximum matching and
start an alternating-path search from an unmatched clause. Every
reached variable must be matched, or an augmenting path would exist.
Its matched clause is then reached as well. The starting unmatched
clause gives more reached clauses than reached variables, which is
the asserted obstruction.

Enlarge a deficient variable set to a set of exactly \(t\) variables.
A union bound gives
\begin{align*}
 \PP(\text{no matching})
 &\le\sum_{t=1}^{m}\binom mt\binom nt(t/n)^{kt}\\
 &\le\sum_{t=1}^{m}(e^2c_0\,t/n)^t=o(1).
\end{align*}
Indeed, the probability that a proper clause lies in a specified
\(t\)-set is \(\fall tk/\fall nk\le(t/n)^k\), and it is zero for
\(t<k\). Bounding the two binomial coefficients by \( (em/t)^t\) and
\((en/t)^t\) first gives \([e^2c_0(t/n)^{k-2}]^t\).
Since \(k\ge3\) and \(t<n\), this is at most the displayed summand.
For the last limit, every fixed summand tends to zero, whereas the tail
is bounded by the tail of the convergent geometric series
\(\sum_{t\ge1}(e^2c_0)^t\).

Thus \(P_n(m)=1-o(1)\). Because \(m\le M_n\), the survival identity
\eqref{eq:survival} gives
\[
 \mu_n\ge\frac mn P_n(m)=c_0-o(1).
\]
An eventual lower bound \(a_0>0\) and the asserted upper bound follow.
\end{proof}
